\section{The contamination operator: conventions and computations}\label{si:contam}

\paragraph{Conventions.} The contaminated recursion \meqref{eq:rho} takes the infimum over $q=(1-\rho)p+\rho r$ with $r$ ranging over all distributions on the outcome set $Y^{\bot}$ of the stage. The adversary may place mass on an outcome $y$ that is nominally null, $\Prob(y\mid b,a)=0$, for which the Bayes posterior is undefined; the operator therefore needs an explicit continuation there. The exact module adopts the convention that a nominally null outcome in $Y$ carries continuation value $0$ (the most pessimistic reading) and that the exit outcome $y_\bot$ always carries value $0$; the main-text statements are insensitive to the choice because in the instances used every outcome in the contaminating support is nominally positive. Fixing the nominal posteriors while perturbing the outcome probabilities makes \meqref{eq:rho} a rectangular stress operator in the sense of robust dynamic programming; it is not the Bayesian update under one coherent perturbed joint law, and no such interpretation is claimed.

\paragraph{The one-stage identity.} Equation \meqref{eq:contam} is the minimum of the linear functional $r\mapsto(1-\rho)\sum_yp(y)v_y+\rho\sum_yr(y)v_y$ over the simplex, attained at a vertex. The module verifies it against explicit vertex enumeration on the instances of the main text.

\paragraph{The non-affine example.} One safe state $s$, one action, $T(s\mid s)=\tfr12$, observations $y_s$ (after a safe transition) and $y_\bot$ (after exit), both nominally of probability $\tfr12$ from $\delta_s$. Then $V^\rho_1(\delta_s)=(1-\rho)\bigl(\tfr12\cdot1+\tfr12\cdot0\bigr)+\rho\min\{1,0\}=\tfr{1-\rho}2$ and, inductively, $V^\rho_{k}(\delta_s)=\bigl(\tfr{1-\rho}2\bigr)^{k}$. At $\rho=\tfr12$ and $k=2$ this is $\tfr1{16}$, whereas the affine interpolation of the endpoint values $V^0_2=\tfr14$ and $V^1_2=0$ gives $\tfr18$. The module computes $V^{1/2}_2$ by the recursion itself.

\paragraph{Prior contamination on the delayed class.} In \meqref{eq:priorcontam} the adversary chooses $r\in[0,1]$ after the hold is announced. On the band $1\le z_0<1+\ell$ the hold $u\equiv+1$ saves the $\theta=+1$ branch only, of contaminated mass $(1-\rho)w+\rho r\ge(1-\rho)w$, and the hold $u\equiv-1$ saves the other branch, of mass $(1-\rho)(1-w)+\rho(1-r)\ge(1-\rho)(1-w)$; the controller's best guaranteed value is $(1-\rho)\max(w,1-w)$. The module verifies the closed form by enumerating the two holds against the two extreme contaminations for $\rho\in\{0,\tfr14,\tfr12,\tfr34,1\}$, $w\in\{0,\tfr1{10},\dots,1\}$ and the grid of $(z_0,\ell)$.

\paragraph{Convexity.} At $\rho=0$ the value is convex (Proposition~\mref{prop:alpha}); for $\rho>0$ the inner minimum is a concave operation and convexity is lost in general. Concavity never holds: the two-floor value $\max(b_0,b_1)$ at $\rho=0$ equals $1$ at both vertices and $\tfr12$ at the midpoint.
